\section{Inferencia heurística de la posesión}
\label{sec:inferencia-posesion}

A partir de la salida del procesamiento visual, el sistema dispone en cada frame de jugadores, porteros, árbitros y balón detectados, asociados temporalmente y proyectados sobre el campo. Sobre esta representación se implementa un módulo heurístico para estimar qué jugador y qué equipo tienen la posesión del balón en cada instante.

El método intenta resolver el problema mediante reglas basadas en proximidad, contacto y continuidad temporal. Para cada frame se localiza el balón y se calcula la distancia entre el balón y el punto inferior de la \textit{bounding box} del resto de detecciones. Además, se comprueba si el balón cae dentro de una región ampliada de la caja del jugador, lo que se considera una señal fuerte de posible control.

La decisión no depende únicamente de la distancia instantánea. También se tienen en cuenta señales dinámicas del balón, como reducciones bruscas de velocidad o cambios de dirección, que pueden indicar un toque. Asimismo, se incorpora continuidad temporal: si un jugador ya tenía la posesión en frames anteriores, el sistema tiende a mantenerla mientras el balón permanezca cerca y no exista evidencia clara de recuperación por parte de un rival. Para aceptar un cambio de posesión se exige una señal más fuerte o una confirmación durante varios frames, reduciendo así oscilaciones provocadas por ruido, oclusiones o errores del tracking.

Cuando no se detecta el balón durante varios frames, o cuando no existe ningún contacto consistente, la posesión se libera para evitar arrastrar una asignación antigua sin evidencia visual. La salida de este módulo es, por tanto, una estimación por frame del jugador poseedor, el equipo poseedor y la razón heurística asociada a la decisión.

Durante el desarrollo del proyecto nos pareció de vital importancia el cálculo de esta heurística ya que considerábamos que iba a ser una señal muy útil de cara a la detección de acciones. Sin embargo, dado el método final escogido para la detección de eventos (PathCRF), y teniendo en cuenta que este método utiliza únicamente la trayectoria de los jugadores como señal de entrada para su modelo, finalmente esta señal del poseedor de la posesión no se consume en ninguna etapa posterior del pipeline.

Durante el desarrollo del proyecto, la estimación heurística de la posesión se consideró una señal potencialmente muy relevante para la detección posterior de acciones. La intuición inicial era que conocer qué jugador o equipo tenía el balón podía facilitar la identificación del jugador protagonista de eventos como pases, controles, tiros o recuperaciones.

Sin embargo, la elección final de PathCRF como método de detección de eventos modificó este planteamiento. Este modelo utiliza como entrada principal las trayectorias de los jugadores y no requiere consumir explícitamente una señal previa de posesión estimada por el pipeline. Por este motivo, aunque el módulo de inferencia heurística de la posesión se implementó y resulta útil como herramienta de análisis y visualización, su salida no se utiliza finalmente como entrada directa en las etapas posteriores del sistema.
